\section{Introduction}
\label{sec:intro}

% Paragraph 1: Motivation
Do small, quasi-random access shocks affect career trajectories in selective labor markets? In tournament-style markets---where ranking determines access and access determines the opportunity to earn further ranking---a single exogenous opportunity may propagate through the system. This paper studies lucky losers in professional tennis: players who lose in the final round of qualifying but enter the main draw when another player withdraws. At Grand Slam tournaments, lucky losers are selected by lottery among the top four ranked qualifying losers, creating a natural experiment in which otherwise similar players are randomly granted or denied a main draw opportunity. The setting allows clean measurement of both tournament access and playing ability, enabling a separation that is rarely possible in other labor markets.

% Paragraph 2: Mechanism intuition
The mechanism behind career propagation in this setting is path dependence through the ranking system. Professional tennis uses a rolling 52-week ranking that determines direct entry into tournament main draws. A player who earns ranking points at one event may gain direct entry to subsequent events, generating further opportunities to earn points. This creates a feedback loop: access generates exposure, exposure generates accumulation, and accumulation generates future access. A one-time exogenous shock to ranking points can trigger this cascade---but only if the initial shock is large enough. Players who receive a lucky loser slot but lose in the first round earn minimal ranking points and trigger no cascade. Players who win one or more main draw matches earn enough points to cross direct-entry thresholds at lower-tier events, setting off a sequence of additional opportunities. The ranking system, however, does not adjust for opponent quality: a first-round win at a Grand Slam earns the same points regardless of the opponent's strength. Whether this ranking inflation translates into worse match outcomes---an ``overseeding'' effect---is an empirical question.

% Paragraph 3: Empirical strategy
The empirical strategy has three components. First, Grand Slam lucky loser allocation is random within the top-four pool, providing the primary identification. I estimate immediate effects at the focal event and dynamic outcomes at horizons $h \in \{4, 8, 12, 26, 52\}$ weeks, tracing the trajectory of access and ability from the immediate aftermath through one year. Second, I extend the analysis to non-Grand-Slam events using a control function approach, where the generalized residual is computed from the exact enumeration probability that a player's rank among qualifying losers falls within the top slots---a probability determined by the full combinatorial structure of qualifying match pairings. Third, a match-level tournament performance model estimates whether lucky loser entry changes subsequent win probability conditional on opponent strength, exploiting the Grand Slam lottery for identification.

% Paragraph 4: Main findings
In the full ATP Grand Slam sample ($N = 248$, with 132 treated and 116 controls across 62 events), the event fixed effects specification shows that lottery selection increases the probability of winning a main draw match by 32 percentage points ($p < 0.001$), adds 1.48 matches played ($p < 0.001$), and generates 29.66 ranking points at the focal event ($p < 0.001$). On the WTA tour ($N = 132$, with 54 treated and 78 controls), the immediate effects are similar: a 38 percentage point increase in winning a main draw match and 24.86 additional ranking points. Over the following weeks, ATP ranking points increase by 31.9 at 4 weeks ($p = 0.028$), 38.5 at 12 weeks ($p = 0.007$), and 49.2 at 26 weeks ($p = 0.054$), fading by 52 weeks as the rolling window replaces the initial gain. WTA effects are larger at short horizons---64.4 points at 4 weeks ($p = 0.003$) and 78.7 at 8 weeks ($p < 0.001$)---but dissipate faster, consistent with the denser WTA tournament calendar. Elo ratings---which adjust for opponent quality---are unchanged at every horizon on both tours. A match-level logit on the Grand Slam lottery sample finds no significant effect on match win probability on either tour: $\hat{\delta} = -0.123$ ($p = 0.119$; 8,543 matches from 130 players) on the ATP tour and $\hat{\delta} = +0.034$ ($p = 0.760$; 5,391 matches from 98 players) on the WTA tour. Ranking point gains come entirely from additional participation rather than improved match performance. The non-Grand-Slam control function estimates corroborate these patterns: ATP ranking point gains of 12.2 at 4 weeks ($p = 0.065$) and 28.1 at 26 weeks ($p = 0.009$), with significant gains in tournament access at longer horizons.

% Paragraph 5: Contribution
This paper makes three contributions. First, I provide the first causal estimates of how a within-career opportunity shock affects trajectories in a tournament labor market. The initial conditions literature has established that graduating in a recession \citep{Oreopoulos2012_graduating_recession}, initial job placement \citep{Oyer2006_initial_placement}, and the environment in which a person starts \citep{Chetty2014_where_is_land} shape long-run outcomes. I test whether these forces operate for a marginal within-career shock in a setting where treatment is randomized and where both access and ability can be measured separately. Second, the access-versus-ability decomposition provides direct evidence on tournament theory: LL entry generates ranking point gains through additional participation, but match-level performance conditional on opponent strength is unchanged on both tours ($\hat{\delta} = -0.123$, $p = 0.119$ on the ATP tour; $\hat{\delta} = +0.034$, $p = 0.760$ on the WTA tour). This cleanly separates the access channel from the ability channel and shows that access alone drives the career effects. Third, \citet{Maity2025_early_career_sports} study lucky losers using correlational methods and find positive associations between entry and career outcomes. I provide the first causal estimates using literal randomization (the Grand Slam lottery), confirming that LL entry generates real access benefits while leaving match-level ability unaffected.

% Paragraph 6: Roadmap
The remainder of the paper proceeds as follows. Section~\ref{sec:literature} reviews the related literature. Section~\ref{sec:data} describes the institutional setting and data. Section~\ref{sec:strategy} presents the econometric framework and identification strategies. Section~\ref{sec:results} reports the main results. Section~\ref{sec:mechanisms} discusses mechanisms. Section~\ref{sec:robustness} presents robustness checks. Section~\ref{sec:conclusion} concludes.
